\section*{Chapter \ref{ch:b2:mixed-strategies} --- Mixed Strategies}

\begin{solution}{exo:b2:mixed-strategies:1}
% code: exercises/solutions.py battle
The pure equilibria are $(A, A)$ and $(B, B)$. In the mixed one, the column player chooses $A$ with probability $q$ that
makes the row player indifferent: $2q = 1 - q$, so $q = 1/3$; the row player chooses $A$ with probability $p$ that makes
the column player indifferent: $p = 2(1 - p)$, so $p = 2/3$. Each player's expected payoff is $2/3$, less than in
either pure equilibrium: the players often fail to coordinate.
\end{solution}

\begin{solution}{exo:b2:mixed-strategies:2}
% code: exercises/solutions.py serve_game
Receiver leans left with $q$: $0.58q + 0.79(1 - q) = 0.73q + 0.49(1 - q)$, so $0.30 = 0.45q$ and $q = 0.667$. Server serves
left with $p$: $0.58p + 0.73(1 - p) = 0.79p + 0.49(1 - p)$, so $0.24 = 0.45p$ and $p = 0.533$. The server wins with
probability $0.58 \times 0.667 + 0.79 \times 0.333 = 0.65$.
\end{solution}

\begin{solution}{exo:b2:mixed-strategies:3}
% figdata: serves_summary
% ledger: b2:ww2001-totals
64 per cent after serves to the left (1{,}040 of 1{,}622) and 65 per cent after serves to the right (918 of 1{,}404).
Equality was rejected in one point game of 40 at the 5 per cent level; if the model is right, each test rejects
with probability 0.05, so about two rejections are expected.
\end{solution}

\begin{solution}{exo:b2:mixed-strategies:4}
Win rates that differ between the two directions systematically, in many point games of the same server, by
more than chance allows. A server who is simply better than another wins more on both sides; that is not
evidence against the model, which concerns the comparison between the sides for one server, against one
receiver, on one court.
\end{solution}

\begin{solution}{exo:b2:mixed-strategies:5}
Row 1 earns more than row 2 against both columns, so the row player plays it; against row 1 the column player
chooses column 2, holding the row player to 2. The pair (row 1, column 2) is a saddle point: 2 is the least
payoff in its row and the greatest in its column. The \omterm{def:b2:mixed-strategies:value}{value} is 2, and no mixing is needed.
\end{solution}

\begin{solution}{exo:b2:mixed-strategies:6}
% code: exercises/solutions.py better_goalie
The keeper's improvement makes kicking left less attractive. Recomputing with 0.50: the keeper covers left with
$q$ where $0.50q + 0.95(1 - q) = 0.90q + 0.70(1 - q)$, so $q = 0.25/0.65 = 0.385$; the kicker kicks left with $p$ where $0.50p +
0.90(1 - p) = 0.95p + 0.70(1 - p)$, so $p = 0.20/0.65 = 0.308$. The \omterm{def:b2:mixed-strategies:value}{value} falls from 0.791 to 0.777. Both kick left and
cover left become rarer: the keeper, now better on the left, covers it \emph{less} often, because the kicker
goes there less and the keeper must keep the kicker indifferent.
\end{solution}

\begin{solution}{exo:b2:mixed-strategies:7}
% code: exercises/solutions.py pearson
Win rates $20/33 = 0.606$ and $22/26 = 0.846$; pooled $42/59 = 0.712$. Left: $(20 - 33 \times 0.712)^2/(33 \times 0.712 \times
0.288) = 1.80$; right: $(22 - 26 \times 0.712)^2/(26 \times 0.712 \times 0.288) = 2.29$; total 4.09. It exceeds 3.84: equality is
rejected at 5 per cent, the one rejection among the 40 point games.
\end{solution}

\begin{solution}{exo:b2:mixed-strategies:8}
If the opponent plays rock, paper and scissors with probabilities $r$, $s$, $c$, rock earns $c - s$, paper $r - c$, scissors
$s - r$. A mixture using all three requires the three to be equal, so $r = s = c = 1/3$, and each earns 0. No pure
strategy is a \omterm{def:b2:static-games:best-response}{best response} to itself and no mixture on two strategies can be part of an equilibrium (the
third strategy would then beat it), so this is the only equilibrium.
\end{solution}

\begin{solution}{exo:b2:mixed-strategies:9}
Count the runs, the unbroken sequences of serves in one direction, in each point game, and compare the
number with its distribution under independent choices with the observed proportions. Too many runs means
the server switches too often; too few, that he repeats. A server who tries to look unpredictable may switch
too often, because sequences that really are random contain longer runs than people expect.
\end{solution}

\begin{solution}{exo:b2:mixed-strategies:10}
Without a saddle point the indifference conditions give $q^* = (a_{22} - a_{12})/d$ and $p^* = (a_{22} - a_{21})/d$ with $d =
a_{11} - a_{12} - a_{21} + a_{22} \ne 0$. The row player's payoff against $q^*$ is $a_{11}q^* + a_{12}(1 - q^*) = (a_{11}a_{22} -
a_{12}a_{21})/d = v$ for both rows. So with $q^*$ the column player holds her to $v$ ($\min\max \le v$), and with $p^*$ she
earns $v$ against both columns ($\max\min \ge v$); with $\max\min \le \min\max$ the three give equality.
\end{solution}

\begin{solution}{exo:b2:mixed-strategies:11}
Let $p$ and $q$ be the probabilities of the first strategies. The row player's \omterm{def:b2:static-games:best-response}{best response} to $q$ depends on the
sign of a linear function of $q$: it is $p = 1$ where the function is positive, $p = 0$ where it is negative, and
every $p$ where it is zero. Its graph is a connected path from the side $q = 0$ to the side $q = 1$ of the unit
square. The column player's is a connected path from $p = 0$ to $p = 1$. Two such paths must cross, and a crossing
point is a pair of mutual \omterm{def:b2:static-games:best-response}{best responses}.
\end{solution}

\begin{solution}{exo:b2:mixed-strategies:12}
% code: exercises/solutions.py inspection
The taxpayer's indifference is unchanged: $a^* = D/(D + F)$. The inspector now earns $D - k + e\,sF$ from auditing and
$(1 - e)D$ from not auditing, so he is indifferent when $e(D + sF) = k$: $e^* = k/(D + sF)$. Evasion now falls with the
fine and with the share $s$, and rises with the cost of an audit; the audit rate still falls with the fine.
\end{solution}

\begin{solution}{pb:b2:mixed-strategies:1}
% code: exercises/solutions.py inspection
\textbf{Part I.}
1. Taxpayer first, office second: honest and audit, $-1{,}000$ and 800; honest and no audit, $-1{,}000$ and
$1{,}000$; evade and audit, $-2{,}500$ and 800; evade and no audit, 0 and 0.
2. Audit $\to$ honest is best $\to$ no audit is best $\to$ evade is best $\to$ audit is best: every cell has a player who
gains by switching.
3. The taxpayer's: the audit rate must make her indifferent.
4. The office's.
5. No: evading is better without an audit, worse with one.
6. No: the taxpayer's payment and the office's net collection do not sum to a constant, because audits cost
and fines go elsewhere.
\textbf{Part II.}
7. $a^* = 1{,}000/2{,}500 = 0.4$.
8. $e^* = 200/1{,}000 = 0.2$.
9. $1{,}000$: honest, she pays it for sure; evading, she pays $0.4 \times 2{,}500 = 1{,}000$ on average.
10. 800: the payoff from auditing, $1{,}000 - 200$, equal to that from not auditing, $0.8 \times 1{,}000$.
11. Tax from honest taxpayers $0.8 \times 1{,}000 = 800$ plus tax recovered from caught evaders $0.2 \times 0.4 \times 1{,}000 =
80$: 880, or 88 per cent; audits cost $0.4 \times 200 = 80$.
12. At $a = 0.4$ the taxpayer pays 1{,}000 either way; at $e = 0.2$ the office collects 800 either way.
\textbf{Part III.}
13. In the model, audits fall and evasion is unchanged.
14. $a^* = 1{,}000/5{,}000 = 0.2$; $e^* = 0.2$.
15. They halve, from 80 to 40 per taxpayer.
16. A cheaper audit, or a larger gain to the office from catching an evader (\cref{exo:b2:mixed-strategies:12}).
17. They may not know the audit rate and judge it from the past, or the fine may carry a stigma the payoffs leave
out; then a higher fine deters before audits adjust.
18. Evasion and audit rates before and after a change in fines, over a period long enough for audits to adjust:
the model predicts audits fall and evasion returns to its old level.
\textbf{Part IV.}
19. $a = 1{,}000/2{,}500 = 0.4$: then evading costs on average as much as paying.
20. $1{,}000 - 0.4 \times 200 - 100 = 820$.
21. $820 > 800$: yes.
22. $a = 1{,}000/1{,}500 = 0.667$, and $1{,}000 - 133.33 - 100 = 766.67 < 800$: no.
23. Commitment pays when $1{,}000 - 200 \times 1{,}000/(1{,}000 + F) - 100 \ge 800$, that is, $200{,}000/(1{,}000 + F) \le 100$, or $F \ge
1{,}000$.
24. In equilibrium the office audits 40 per cent of taxpayers and 20 per cent evade; announcing an audit rate pays
when the fine is at least 1{,}000 dalens.
\end{solution}
